\section{Limits of Current Defences}
\label{sec:limits}

The defences covered in the previous sections represent meaningful progress,
but the research also makes clear that none of them fully solve the problem.
This section examines why --- both in terms of how defences are evaluated and
what happens when attackers target the defence mechanisms directly.

\subsection{Evaluation Methodology Problems}

A recurring issue is that reported robustness often does not reflect realistic
adversarial conditions. \textcite{jia2025critical} argue that many published
evaluations rely on assumptions too narrow to support their conclusions.

The first problem is \textit{attack diversity}. Most evaluations test defences
against heuristic attacks --- handwritten prompts attempting to override system
instructions. These provide a useful baseline, but they do not represent an
attacker who knows what defence is in place and targets it
specifically~\parencite{jia2025critical}.

The second is \textit{prompt diversity}. StruQ and SecAlign both reported
strong results in their original evaluations, but those tests used a limited
set of injection prompts. When evaluated against a broader range, success rates
increase considerably~\parencite{jia2025critical}. The defence does not
perform worse --- the evaluation is simply more representative.

Finally, \textit{utility preservation} is difficult to measure consistently.
Different benchmarks yield substantially different conclusions about how much
capability a defence costs. A system can appear secure and functional under one
evaluation and noticeably degraded under another~\parencite{jia2025critical},
which makes fair comparison between defences difficult.

\subsection{Adaptive Attacks}

The clearest evidence that current defences are less robust than reported comes
from adaptive attacks --- attacks designed specifically to target the defence
mechanism rather than the model in general.

\textcite{yang2025checkpointgcg} introduce Checkpoint-GCG, which exploits
intermediate fine-tuning checkpoints to construct adversarial suffixes that
transfer to the fully trained model. The results show that SecAlign's
robustness against standard GCG does not imply the underlying vulnerability
has been resolved --- the attack surface remains, but requires a more targeted
approach to reach~\parencite{yang2025checkpointgcg}.

\textcite{pandya2025astra} take a different approach with ASTRA, which targets
the model's attention mechanism rather than output token probabilities. Both
ASTRA and its universal variant ASTRA++ achieve high success rates even under
weak-knowledge assumptions, demonstrating that fine-tuning alone is not
sufficient to reliably prevent prompt injection~\parencite{pandya2025astra}.
These results also raise questions about GCG as an evaluation tool: because
gradient-based optimisation provides a weak signal against fine-tuned models,
poor GCG performance does not necessarily indicate strong
robustness~\parencite{pandya2025astra}.

Taken together, these results suggest that the fundamental weakness of
fine-tuning-based defences is not the failure of any individual method, but
the fact that behavioural alignment itself provides no hard separation between
instructions and data. Improvements in robustness therefore appear to increase
the effort required for attacks rather than eliminate the underlying attack surface.

\subsection{Security and Utility Trade-Offs}

Every defence discussed in this paper introduces a trade-off between security
and utility. Detection-based approaches reduce attack success rates but
generate false positives that interfere with legitimate use
cases~\parencite{debenedetti2024agentdojo, jia2025critical}. Fine-tuning
approaches impose measurable capability losses~\parencite{jia2025critical}.
The underlying reason is that robustness and instruction-following are in
tension: a model trained to treat data-channel instructions with suspicion will
sometimes apply that suspicion to legitimate inputs as well. SecAlign handles
this better than StruQ, but the trade-off cannot be fully
eliminated~\parencite{chen2025secalign}.

\subsection{Why Prompt Injection Remains an Open Problem}

Across the research reviewed here, the same pattern recurs: defences that
appear robust under standard evaluation conditions weaken significantly when
tested against adaptive attackers or more comprehensive benchmark
sets~\parencite{jia2025critical, yang2025checkpointgcg}.

The underlying reason is structural. Prompt injection exploits the core
capability of LLMs --- following natural language instructions. As long as
instructions and external data share the same context window, that attack
surface exists. Fine-tuning increases the cost of a successful attack, but
does not eliminate the condition that makes it possible.

No currently available defence provides complete protection against prompt
injection. Understanding it as a system-level problem rather than a model-level
one has direct implications for secure deployment, which are discussed in the
following section.
